\documentclass[b5j]{jsarticle}
% このファイルは pLaTeX でコンパイルします。
% 例: platex aQUIZZ.tex && dvipdfmx aQUIZZ.dvi

\usepackage{amsmath,amssymb,mathtools}
\usepackage{enumitem}
\usepackage{xcolor}
\usepackage[dvipdfmx]{graphicx}
\usepackage[dvipdfmx]{hyperref}
\usepackage{fancyhdr}
\newcommand{\bou}[2]{%
\underline{#2}\fbox{\scriptsize #1}
}
\usepackage[
  truedimen,
  b5j,
  left=15mm,
  right=15mm,
  top=12mm,
  bottom=12mm
]{geometry}

\pagestyle{fancy}
\fancyhf{}
\fancyhead[C]{第二回　恋愛共通テスト}
\fancyfoot[C]{\thepage}

\setlength{\parindent}{1zw}
\setlength{\parskip}{0.3em}

\begin{document}

%-------------------------
% 表紙部分
%-------------------------

\begin{center}

\vspace*{10em}

{\Huge\bfseries 第二回　恋愛共通テスト}

\vspace{1em}

{\Large 恋愛科目}

\vspace{2em}

{\large 2026年11月実施}

\end{center}

\vspace{2em}



\vspace{10em}

\begin{center}
\fbox{
\begin{minipage}{0.9\linewidth}
\small
\textbf{受験上の注意}

\begin{enumerate}[leftmargin=2em]
\item 問題冊子は監督者の指示があるまで開いてはいけない。
\item 解答に際しては、論理的思考力・状況判断力・察する力を総合的に活用すること。
\item 出題内容は実際の恋愛状況と異なる場合がある。
\item 必要に応じて常識・経験・反省を用いて解答すること。
\item 本問の採点はあくまでも採点者の主観であり、正解は存在しない。
\item 解答は司会により示されたタイミングでフリップに書いて示し、カメラに見えるようにしながら読み上げること。 
\end{enumerate}

\end{minipage}
}
\end{center}
\vspace{2em}
\begin{flushright}
75期ステージ班
\end{flushright}
\newpage

%-------------------------
% 第一問
%-------------------------

\begin{center}
{\Large\bfseries 第一問〔必答問題〕}

\vspace{0.5em}

（配点50点）
\end{center}

\vspace{1em}

\noindent
〔1〕次の会話を読み、後の問いに答えよ。

\vspace{0.5em}

\begin{quote}
  （AさんとBさんはデートで買い物にきました。Aさんはお気に入りのコスメを探しています）


A：「うーん、ないなあ。」

B：「このブランドの名前、なんて読むの？」

A：「アイオペっていうんだけど...」

B：「\bou{1}{（アイオペ...?どんなスペルなんだろう）}」

A：「あ、あった！これこれ！この\bou{2}{ティント！}」

\end{quote}

\vspace{1em}

\begin{enumerate}
\item Bさんが知らなかったコスメブランド、アイオペのスペルを答えよ。
\fbox{
  \begin{minipage}[c][1em][c]{1em}
    \centering
    1
  \end{minipage}
}



\vspace{5em}

\item 傍線部2のティントとは何か、リップとの違いを明確にしながら示せ。

\end{enumerate}
〔2〕以下のテーマに関するあなたの意見を述べよ　
\vspace{1em}
\begin{quote}
「あなたの理想の東京の雨デートプラン」
\end{quote}
\newpage
%-------------------------
% 第二問
%-------------------------
\begin{center}
{\Large\bfseries 第二問〔必答問題〕}

\vspace{0.5em}
（配点50点）
\end{center}
\vspace{1em}
\noindent
〔1〕次の問に答えよ
\vspace{0.5em}
\begin{quote}
  ちー
\end{quote}
\newpage
\vspace{10em}
\fbox{
\begin{minipage}{0.9\linewidth}
\small
\begin{center}
{\textbf{【採点基準】}
  \begin{enumerate}[leftmargin=2em]
\item 相手の気持ちを考慮しているか
\item 現実性があるか
\item 興味深いか
  \end{enumerate}

}
\end{center}
\end{minipage}
}
\end{document}